\documentclass[twocolumn]{revtex4-2}

\usepackage{amsmath,amssymb}
\usepackage{graphicx}
\usepackage{booktabs}
\usepackage{xcolor}
\usepackage{hyperref}

\newcommand{\GN}{G_N}
\newcommand{\GNemerg}{G_N^{\mathrm{emerg}}}
\newcommand{\GNS}{G_N^{\mathrm{Sakharov}}}
\newcommand{\Lcut}{\Lambda_{\mathrm{UV}}}
\newcommand{\Nf}{N_f}
\newcommand{\alphaADW}{\alpha_{\mathrm{ADW}}}
\newcommand{\Dslash}{D\!\!\!\!/}

% Auto-generated by update_counts.py — 2026-04-24T19:44:03.140400
% Do not edit manually. Run: uv run python scripts/update_counts.py
\newcommand{\totaltheorems}{3993}
\newcommand{\substantivetheorems}{3882}
\newcommand{\placeholdertheorems}{111}
\newcommand{\axiomcount}{1}
\newcommand{\sorrycount}{0}
\newcommand{\sorrypercent}{0.0}
\newcommand{\leanmodules}{167}
\newcommand{\leandefinitions}{3297}
\newcommand{\pythonmodules}{68}
\newcommand{\testfiles}{59}
\newcommand{\totaltests}{2836}
\newcommand{\figurecount}{111}
\newcommand{\notebookcount}{54}
\newcommand{\papercount}{19}
\newcommand{\aristotleproved}{322}
\newcommand{\aristotleruns}{44}
% --- Paper 18 / Phase 5t per-section counts (FermiHubbardDimer.lean) ---
\newcommand{\fhdTotal}{143}
\newcommand{\fhdWaveTwo}{13}
\newcommand{\fhdWaveThree}{6}
\newcommand{\fhdWaveFour}{22}
\newcommand{\fhdWaveFive}{22}
\newcommand{\fhdWaveSixABC}{38}
\newcommand{\fhdWaveSixExtra}{15}
\newcommand{\fhdWaveSeven}{19}
\newcommand{\fhdWaveEight}{8}

% \heatKernelThms and \heatKernelTests are auto-generated by
% scripts/update_counts.py from the canonical Lean source +
% pytest count and live in docs/counts.tex.

\begin{document}

\title{Formal heat-kernel calibration of induced Newton's constant from
the ADW Dirac substrate}

\author{John G. Roehm}
\affiliation{Independent Researcher; \texttt{NetRxn137@gmail.com}}

\date{\today}

\begin{abstract}
We formalize, in the Lean 4 proof assistant, the leading
Seeley--DeWitt heat-kernel coefficients of the Dirac
operator on a $4$-dimensional Riemannian background and prove that
the coefficient $a_2$ of the Ricci scalar reproduces the Sakharov--Adler
induced Newton constant
$\GNS = 12\pi/(\Nf \Lcut^2)$~\cite{Sakharov1968,Adler1982}
\emph{exactly}, providing a closed-form structural anchor for the
emergent-gravity claim of the ADW (Alvarez--Diakonov--Wetterich)
mean-field framework. The calibration is a load-bearing input for the
non-linear effective action: any deviation of the heat-kernel
$a_2$ coefficient from the linearized
$\GNemerg(\Lcut, \Nf, \alphaADW = 1)$ value of Phase 6a.1 would mark
the mean-field validity boundary. We prove a biconditional that
sharpens this gate: the heat-kernel $\GN$ matches the linearized
$\GNemerg$ at given $(\Lcut, \Nf)$ \emph{if and only if} the
ADW dimensionless coefficient $\alphaADW$ equals~$1$ (the Sakharov
baseline). The Lean module
\texttt{HeatKernelExpansion.lean} ships
\heatKernelThms{} substantive theorems, zero \texttt{sorry} statements,
and zero new axioms beyond the Lean~4 kernel's standard three.
\textbf{The PDE-level asymptotic-expansion theorem itself is not
re-derived in Lean} (it requires Mathlib's spin-bundle
infrastructure, which is not yet available); instead it is
\emph{encoded as a tracked hypothesis} via the structure
\texttt{DiracHeatKernelAsymptotic}, whose invariants force any
downstream consumer to commit to the Christensen--Duff--Vassilevich
textbook coefficient
table~\cite{ChristensenDuff1979,Vassilevich2003}. All
algebraic content downstream of the structure --- coefficient
identities, Decision-Gate-E.2 calibration, Sakharov--Adler match ---
is fully formalized. A Python mirror in
\texttt{src/heat\_kernel/} provides numerical evaluation of the
coefficient triple and a \heatKernelTests{}-case \texttt{pytest} suite
cross-checks the Lean theorems against the Christensen--Duff
rationals.
\end{abstract}

\maketitle

\section{Introduction}
\label{sec:intro}

The proposal that gravity in our world is an effective long-distance
description of a microscopic fermionic substrate dates to
Sakharov's 1968 induced-gravity programme~\cite{Sakharov1968}, refined
by Adler~\cite{Adler1982} and revived in modern form by
Diakonov~\cite{Diakonov2011} and the Alvarez--Diakonov--Wetterich
(ADW) tetrad-condensation framework~\cite{VladimirovDiakonov2012}.
In the Sakharov--Adler limit, the Newton constant is induced by
integrating out $\Nf$ Dirac fermions at one loop with a hard
ultraviolet cutoff $\Lcut$:
\begin{equation}
  \GNS \;=\; \frac{12\pi}{\Nf \, \Lcut^2}.
  \label{eq:gnsakharov}
\end{equation}
Phase~6a.1 of this project formalized the linearized Einstein equations
emerging from the ADW substrate and proved that the leading
emergent Newton constant takes the form
$\GNemerg(\Lcut, \Nf, \alphaADW) = \alphaADW \cdot \GNS$, with
$\alphaADW = 1$ reproducing Eq.~\eqref{eq:gnsakharov} exactly and
$\alphaADW \ne 1$ encoding ADW-specific corrections away from the
free-fermion limit~\cite{LinearizedEFE2026,Vergeles2025}.

Phase~6e steps up to the \emph{non-linear} effective action: the full
Einstein--Hilbert action plus higher-curvature corrections, derived as
a Seeley--DeWitt heat-kernel expansion of the Dirac fermion
determinant
$\log\det(\Dslash) = -\tfrac{1}{2}\int_0^\infty \tfrac{d\tau}{\tau}\,
\mathrm{Tr}\,\exp(-\tau\,\Dslash^{\,2})$. Wave~1 of Phase~6e --- the
present work --- formalizes the leading-order coefficients
$a_0, a_2, a_4$ of this expansion, derives the emergent
Newton constant by matching the Einstein--Hilbert prefactor against
the $\Lcut^2 \cdot a_2$ divergence, and proves that the result agrees
with the linearized 6a.1 prediction \emph{exactly} at
$\alphaADW = 1$.

The match is not a check of an unknown coefficient against a known
one; it is a structural cross-bridge that anchors the non-linear
calculation to the linearized one. Crucially, the same Lean module
that proves the calibration also proves a biconditional: the
heat-kernel $\GN$ agrees with $\GNemerg$ at given $(\Lcut, \Nf)$ if
and only if $\alphaADW = 1$. Departures from $\alphaADW = 1$ produced
by ADW-specific corrections away from the free-fermion limit therefore
appear in the heat-kernel calculation only as departures from the
Sakharov--Adler $a_2$ closed form --- never as ambiguities in the
calibration relation itself. This makes Decision Gate E.2 of the
project's Phase~6e roadmap a clean falsifier: future deep-research
returns on the Vergeles unitarity coefficient $\alphaADW$ are
quantitatively comparable to the heat-kernel rational
coefficient $-1/(12)$ via this biconditional.

The remainder of the paper is organized as follows.
Section~\ref{sec:asymptotic} reviews the heat-kernel asymptotic
expansion and lists the Christensen--Duff--Vassilevich
spin-$\tfrac{1}{2}$ coefficients in the Lichnerowicz convention.
Section~\ref{sec:calibration} presents the formal calibration
theorem $\GN^{(a_2)} = \GNS$ and its quantitative anchor at the
GUT-scale fiducial parameters.
Section~\ref{sec:gate} proves the Decision Gate E.2 biconditional
matching the heat-kernel and linearized Newton constants iff
$\alphaADW = 1$. Section~\ref{sec:a4} records the rational
coefficients of the higher-curvature $a_4$ basis and the local
Gauss--Bonnet algebraic identity. Section~\ref{sec:discussion}
discusses the mean-field validity boundary and what would falsify it.
The Lean module structure and tracked-hypothesis disclosures appear in
Section~\ref{sec:lean}.

\section{Heat-kernel asymptotic and Dirac coefficients}
\label{sec:asymptotic}

For a self-adjoint elliptic operator $\Delta$ of Laplace type
on a smooth Riemannian $4$-manifold, the trace of its heat semigroup
admits the small-$\tau$ asymptotic
expansion~\cite{Gilkey1995,Vassilevich2003}
\begin{equation}
  \mathrm{Tr}\,e^{-\tau \Delta} \;\sim\;
  \frac{1}{(4\pi\tau)^2}\,\sum_{n \ge 0} \tau^n \int d^4 x \sqrt{g}\,
  a_{2n}(x), \quad \tau \to 0^+,
  \label{eq:hkexp}
\end{equation}
with universal coefficient densities $a_{2n}(x)$ that are local
polynomials in the curvature, the gauge field, and the endomorphism
$E$ entering $\Delta = -g^{\mu\nu} \nabla_\mu \nabla_\nu - E$.

For the $4$-dimensional free Dirac operator $\Dslash = \gamma^\mu
\nabla_\mu$ in vacuum (no torsion, no gauge field), the relevant
Lichnerowicz square is $\Dslash^{\,2} = -\nabla^2 + R/4$, so
$E = -R/4$. The standard textbook coefficients per Dirac species,
in the Christensen--Duff convention, are
\begin{align}
  a_0(x) &= \mathrm{tr}\,\mathbf{1}_4 = 4,
  \label{eq:a0}\\
  a_2(x) &= -\frac{R(x)}{12} \cdot \mathrm{tr}\,\mathbf{1}_4 = -\frac{R(x)}{3},
  \label{eq:a2}\\
  a_4(x) &= -\frac{1}{180}\Big[
    \tfrac{5}{12} R^2 - \tfrac{7}{12} R_{\mu\nu} R^{\mu\nu}
    + R_{\mu\nu\rho\sigma} R^{\mu\nu\rho\sigma}\Big],
  \label{eq:a4}
\end{align}
in units where the prefactor $1/(4\pi)^2$ from
Eq.~\eqref{eq:hkexp} is left implicit. The Lean module factors out
the $(4\pi)^2$ Gaussian normalization explicitly via a
\texttt{noncomputable} constant
\texttt{fourPiSqInv = 1/(4$\pi$)$^2$}, with positivity proved by
composing $\mathrm{Real.pi\_pos}$ and \texttt{pow\_pos}. The species
multiplicity factor $\Nf$ enters linearly through the per-species
trace.

The PDE-level statement of Eq.~\eqref{eq:hkexp} requires
manifold-and-spin-bundle infrastructure not yet available in
Mathlib. The Lean module therefore encodes the asymptotic content as
a tracked-hypothesis structure
\texttt{DiracHeatKernelAsymptotic} whose invariants pin the
coefficients to Eqs.~\eqref{eq:a0}--\eqref{eq:a2} via a definitional
equality
$\texttt{a2\_R\_value} : \texttt{a2\_R\_coef} = -(N_f/12)/(4\pi)^2$. Any
downstream consumer of the structure thereby commits to the
Christensen--Duff coefficient table without re-proving the
Gilkey--Vassilevich theory in line.

\section{Calibration to Sakharov--Adler}
\label{sec:calibration}

The Einstein--Hilbert action $S_{\mathrm{EH}} = -(1/16\pi\GN) \int R
\sqrt{g} \, d^4 x$ enters the effective action of the
proper-time-regularized fermion determinant by integrating the
$\Lcut^2$-divergent piece of $a_2$ in Eq.~\eqref{eq:hkexp}
against $\sqrt{g}\,d^4x$. The resulting integral is
\begin{equation}
  -\frac{1}{16\pi\GN}\int R \sqrt{g}\,d^4x
  \;=\; \frac{1}{(4\pi)^2}
        \int d^4 x \sqrt{g}\, \Lcut^2 \, a_2(x)
        \, \Nf,
  \label{eq:eh-match}
\end{equation}
which fixes the Newton constant to
\begin{equation}
  \GN \;=\; \frac{12\pi}{\Nf \Lcut^2}\,,
  \label{eq:gnhk}
\end{equation}
in exact agreement with Eq.~\eqref{eq:gnsakharov}.

The Lean module formalizes Eq.~\eqref{eq:gnhk} via the definition
\texttt{G\_N\_from\_a2 := 12$\pi$ / ($N_f \Lambda^2$)} together with
the substantive cross-bridge
$\texttt{G\_N\_from\_a2\_eq\_G\_N\_sakharov}$
whose proof body invokes
\texttt{LinearizedEFE.G\_N\_sakharov} by name. This protects against
documentation drift: the formula docstring referencing the
6a.1 module is backed by an actual Lean call rather than a
free-text reference.

A quantitative anchor at the GUT-scale fiducial parameters
$(\Lcut, \Nf) = (10^{16}\,\mathrm{GeV}, 15)$ ($\Lcut, \Nf > 0$
throughout) closes the form:
\begin{equation}
  \frac{1}{\GN^{(a_2)}(\Lcut = 10^{16}, \Nf = 15)}
  \;=\; \frac{15 \cdot 10^{32}}{12\pi}
  \;\approx\; 3.98\times 10^{31}\,\mathrm{GeV}^2,
  \label{eq:gut-anchor}
\end{equation}
which sits an order of magnitude below the Planck-mass square
$M_P^2 \approx 1.49\times 10^{38}\,\mathrm{GeV}^2$, formalized in
Lean as the strict-inequality theorem
\texttt{G\_N\_from\_a2\_inverse\_at\_GUT\_below\_planck\_squared}
proved via \texttt{Real.pi\_gt\_three} plus \texttt{norm\_num}. The
inverse-form anchor
\texttt{G\_N\_from\_a2\_at\_GUT\_inverse} clears the nested division
$1/(a/b) = b/a$ via \texttt{one\_div\_div}, exposing the EH prefactor
directly.

\section{Decision Gate E.2 biconditional}
\label{sec:gate}

The match in Section~\ref{sec:calibration} is the linchpin
correctness-push anchor. We sharpen it to a biconditional:

\begin{quote}
\textbf{Theorem (Decision Gate E.2).} \emph{Let $\Lcut, \Nf > 0$ and
$\alpha \in \mathbb{R}$. Then}
\begin{equation*}
  \GN^{(a_2)}(\Lcut, \Nf) = \GNemerg(\Lcut, \Nf, \alpha)
  \quad \Longleftrightarrow\quad \alpha = 1.
\end{equation*}
\end{quote}

\noindent\emph{Proof.} The forward direction substitutes
Eq.~\eqref{eq:gnhk} for the LHS, expands
$\GNemerg = \alpha\,\GNS$ for the RHS, and uses
$\GNS \ne 0$ (proved via
\texttt{LinearizedEFE.G\_N\_sakharov\_pos}) to right-cancel and
conclude $\alpha = 1$ via \texttt{mul\_right\_cancel}. The reverse
direction substitutes $\alpha = 1$ and dispatches to
\texttt{LinearizedEFE.G\_N\_emerg\_at\_alpha\_one}, then to the
calibration identity. Both directions are formalized in
\texttt{a2\_matches\_GNemerg\_iff\_alpha\_ADW\_unity}.

The biconditional is not a definitional restatement: the LHS is
strictly stronger than $\alpha = 1$ because it specifies that the
match holds at given $(\Lcut, \Nf)$, and the proof of the forward
direction explicitly invokes the non-vanishing of $\GNS$ on the
positive cone $\Lcut, \Nf > 0$. This is the formal sense in which
the heat-kernel calibration constrains $\alphaADW$: any future
deep-research return that produces $\alphaADW \ne 1$ from
microscopic ADW unitarity arguments
falsifies the simultaneous hypotheses (a) ``standard
Christensen--Duff $a_2$ coefficient'' and (b) ``Sakharov--Adler
linearized Newton constant.'' The biconditional makes the deep
research's answer either compatible (if it returns $\alphaADW = 1$) or
diagnostic (if it does not).

\begin{figure}[t]
  \centering
  \includegraphics[width=\columnwidth]{figures/fig_a2_vs_linearized_G_N.png}
  \caption{Phase~6e Wave~1 calibration. Left: relative error
  $|G_N^{(a_2)} - \GNemerg|/\GNemerg$ over $\alphaADW \in [0.05, 5.0]$
  at GUT-anchor parameters $(\Lcut, \Nf) = (10^{16}, 15)$; exact zero
  at $\alphaADW = 1$, with the project's $\pm 50\%$ pass band
  (matching \texttt{GRAV\_PARAMS.G\_N\_MATCH\_TOLERANCE}) shaded.
  Right: log-log $\GN^{(a_2)}(\Lcut)$ at fixed $\Nf = 15$ over
  $[10^{10}, 10^{19}]\,\mathrm{GeV}$ vs the CODATA reference
  $G_N^{\mathrm{obs}}$; the curve crosses the observed value near
  the Planck mass anchor, the natural scale at which the
  Sakharov--Adler induced Newton constant matches observation.}
  \label{fig:cal}
\end{figure}

\section{Higher-curvature $a_4$ basis}
\label{sec:a4}

The $a_4$ density of Eq.~\eqref{eq:a4} is a local polynomial in three
curvature scalars: $R^2$, $R_{\mu\nu} R^{\mu\nu}$,
$R_{\mu\nu\rho\sigma} R^{\mu\nu\rho\sigma}$. The Christensen--Duff
rationals fix the Dirac contributions:
\begin{align}
  c_{R^2}^{\mathrm{Dirac}} &= -\frac{5}{12 \cdot 180}\,\frac{\Nf}{(4\pi)^2},
  \label{eq:cR}\\
  c_{R_{\mu\nu}^2}^{\mathrm{Dirac}} &= +\frac{7}{12 \cdot 180}\,\frac{\Nf}{(4\pi)^2},
  \label{eq:cRicci}\\
  c_{R_{\mu\nu\rho\sigma}^2}^{\mathrm{Dirac}} &= -\frac{12}{12 \cdot 180}\,\frac{\Nf}{(4\pi)^2}.
  \label{eq:cRiem}
\end{align}
The Lean module proves the sign theorems
\texttt{a4\_R\_sq\_coef\_neg}, \texttt{a4\_Ricci\_sq\_coef\_pos},
\texttt{a4\_Riemann\_sq\_coef\_neg} via \texttt{nlinarith} composing
\texttt{fourPiSqInv\_pos} with the rational coefficients.

The Gauss--Bonnet density
$\mathcal{G} = R^2 - 4 R_{\mu\nu} R^{\mu\nu} +
R_{\mu\nu\rho\sigma} R^{\mu\nu\rho\sigma}$ is topological in 4D; on a
closed manifold $\int_M \mathcal{G}\,\sqrt{g}\,d^4 x = 32\pi^2\chi(M)$.
At the local algebra level, the linear combination of
Eqs.~\eqref{eq:cR}--\eqref{eq:cRiem} evaluates to
\begin{equation}
  c_{R^2}^{\mathrm{Dirac}} - 4 c_{R_{\mu\nu}^2}^{\mathrm{Dirac}}
  + c_{R_{\mu\nu\rho\sigma}^2}^{\mathrm{Dirac}}
  \;=\; -\frac{45}{12 \cdot 180}\,\frac{\Nf}{(4\pi)^2}
  \;=\; -\frac{\Nf}{48\,(4\pi)^2}.
  \label{eq:gb-local}
\end{equation}
This is non-zero locally; the topological vanishing arises only after
integrating against the volume measure on a closed
$M$. The Lean theorem
\texttt{a4\_gauss\_bonnet\_combination} formalizes
Eq.~\eqref{eq:gb-local} via \texttt{ring}, providing a
local-algebra sanity check that the rational coefficient triple is
consistent. Manifold-level integration is deferred to Phase~6f.1
\texttt{Curvature.lean}.

\section{Discussion: mean-field validity}
\label{sec:discussion}

Decision Gate E.2 of the Phase~6e roadmap defines the mean-field
validity boundary as the $\pm 50\%$ band around $\alphaADW = 1$ in
the natural-parameter region $(0.5 \le \alphaADW \le 1.5)$.\footnote{%
This is a tighter band than the broader Vergeles natural range
$\alphaADW \in [0.1, 10]$ used in companion paper~23
(\texttt{paper23\_linearized\_efe})~\cite{LinearizedEFE2026}: the
broader range encodes Vergeles-unitarity admissibility, while the
tighter ($\pm 50\%$) band here encodes the heat-kernel-vs-linearized
calibration agreement at the
\texttt{GRAV\_PARAMS.G\_N\_MATCH\_TOLERANCE} value. Both coexist:
$[0.5, 1.5]$ is the calibration-match window inside the broader
unitarity-admissibility region $[0.1, 10]$.}\ Inside
the band, the heat-kernel calibration agrees with the linearized
$\GNemerg$ to better than 50\% relative error, and ADW emergent
gravity is mean-field-consistent at order $\Lcut^2$. Outside the
band (e.g., a Vergeles unitarity return giving $\alphaADW = 5$), the
mean-field approximation is provably inadequate at the $a_2$
calibration --- a falsifiable consequence of the formal biconditional.

The biconditional sharpness is tighter than the validity-band claim:
\emph{at} $\alphaADW = 1$ the match is exact; \emph{anywhere else} the
match has positive relative error. The $\pm 50\%$ band is a project
policy choice (matching the GRAV\_PARAMS tolerance), not a Lean
theorem; but the Lean theorem
\texttt{a2\_matches\_GNemerg\_iff\_alpha\_ADW\_unity} forces any
future tolerance refinement to be expressible as a strict-inequality
constraint on $\alphaADW$ alone, since
$(\Lcut, \Nf)$ factor out completely.

This sharpness has a downstream consequence for the rest of Phase~6e.
Wave~5 (microscopic-to-macroscopic coefficient match) propagates the
$\alphaADW$ value into the cosmological-constant scale $\Lambda^{\mathrm{emerg}}$
through the $a_0$ coefficient. If $\alphaADW = 1$ from the
present formalization is consistent with Vergeles 2025 unitarity,
then $\Lambda^{\mathrm{emerg}}$ inherits the Sakharov--Adler structural form
and the cosmological-constant problem is reduced to fixing the
microscopic UV cutoff $\Lcut$. If $\alphaADW \ne 1$ from the
deep-research return, the present biconditional flags the discrepancy
as a non-perturbative ADW correction with a clean falsifier at
$a_2$; the heat-kernel coefficient becomes a probe of the
mean-field validity rather than a cross-check on the linearized
calculation.

\section{Lean module structure}
\label{sec:lean}

The Lean module \texttt{lean/SKEFTHawking/HeatKernelExpansion.lean}
ships \heatKernelThms{} substantive theorems organized in eight
sections matching the structure of this paper:
\S 1 (Gaussian normalization, 2 theorems);
\S 2 (leading $a_0$ coefficient, 2 theorems);
\S 3 (Einstein--Hilbert $a_2$ coefficient, 3 theorems including
the vanishing biconditional
\texttt{a2\_R\_coefficient\_eq\_zero\_iff});
\S 4 (tracked structure
\texttt{DiracHeatKernelAsymptotic}, 3 theorems consuming structure
invariants);
\S 5 (Sakharov--Adler calibration, 2 theorems including the
load-bearing cross-bridge
\texttt{G\_N\_from\_a2\_eq\_G\_N\_sakharov});
\S 6 (quantitative anchors, 2 theorems including the
\texttt{norm\_num} GUT-vs-Planck-squared falsifier);
\S 7 ($a_4$ higher-curvature basis, 4 theorems including
\texttt{a4\_gauss\_bonnet\_combination}); and
\S 8 (the Decision Gate E.2 biconditional
\texttt{a2\_matches\_GNemerg\_iff\_alpha\_ADW\_unity}).

The module ships zero \texttt{sorry} statements and zero new axioms
beyond the Lean~4 kernel's standard
\texttt{propext}, \texttt{Classical.choice}, \texttt{Quot.sound}.
The asymptotic-expansion existence (Eq.~\eqref{eq:hkexp}) is encoded
as a tracked hypothesis via the
\texttt{DiracHeatKernelAsymptotic} structure: any user instantiating
the structure commits to the Christensen--Duff--Vassilevich
coefficient values via the structure's
\texttt{a0\_value} and \texttt{a2\_R\_value} invariants. The
manifold-level statement requires Mathlib's spin-bundle infrastructure
which is not yet available; the structure formalizes the algebraic
content without re-deriving the PDE asymptotics.

A Python mirror in \texttt{src/heat\_kernel/} provides numerical
evaluators for $a_0, a_2, a_4$ via
\texttt{seeley\_dewitt\_coefficients}, the calibration check via
\texttt{a2\_calibration\_passes}, and a
\heatKernelTests{}-case \texttt{pytest} suite cross-checking each
Lean theorem against the Christensen--Duff rational
table. Canonical formulas live in \texttt{src/core/formulas.py} as
\texttt{seeley\_dewitt\_a0}, \texttt{seeley\_dewitt\_a2\_R\_coefficient},
\texttt{seeley\_dewitt\_a4\_basis}, \texttt{G\_N\_from\_seeley\_dewitt},
and \texttt{heat\_kernel\_a2\_matches\_GN\_sakharov}, each with explicit
\texttt{Lean:} cross-references to the corresponding theorem.

\section{Conclusion}
\label{sec:conclusion}

The Sakharov--Adler induced-gravity coefficient
$\GNS = 12\pi/(\Nf \Lcut^2)$ is recovered \emph{exactly} from the
heat-kernel $a_2$ coefficient of the free Dirac operator. The
formalized result is sharper than a numerical match: it is a
biconditional in which any deviation from the textbook Christensen--Duff
rational
$-1/12$ is detectable via the ADW dimensionless coefficient
$\alphaADW \ne 1$ in the linearized 6a.1 calculation, and the mismatch
factors entirely through the calibration relation rather than through
the cutoff or species count. This makes the heat-kernel $a_2$
calibration a falsifiable mean-field consistency anchor for the rest of
Phase~6e: subsequent waves (higher-curvature structure, non-linear
diff-invariance, full Einstein equations, microscopic-to-macroscopic
coefficient match, Einstein--Cartan extension) inherit a load-bearing
cross-bridge to the linearized 6a.1 module that this Wave~1 closes
with explicit Lean theorems. Whether the bridge survives the Vergeles
unitarity computation will fix the
mean-field validity boundary in the natural-parameter band; either
outcome is substantive.

\begin{thebibliography}{99}

\bibitem{Sakharov1968}
A.~D.~Sakharov,
\emph{Vacuum quantum fluctuations in curved space and the theory of
gravitation},
Sov. Phys. Dokl. \textbf{12}, 1040 (1968);
reprinted in Gen. Rel. Grav. \textbf{32}, 365 (2000).

\bibitem{Adler1982}
S.~L.~Adler,
\emph{Einstein gravity as a symmetry-breaking effect in quantum field
theory},
Rev. Mod. Phys. \textbf{54}, 729 (1982); erratum
Rev. Mod. Phys. \textbf{55}, 837 (1983).

\bibitem{Diakonov2011}
D.~Diakonov,
\emph{Towards lattice-regularized quantum gravity},
arXiv:1109.0091 (2011).

\bibitem{VladimirovDiakonov2012}
A.~A.~Vladimirov and D.~Diakonov,
\emph{Phase transitions in spinor quantum gravity on a lattice},
Phys. Rev. D \textbf{86}, 104019 (2012);
arXiv:1208.1254.

\bibitem{Vergeles2025}
S.~N.~Vergeles,
\emph{Unitarity of 4D Lattice Theory of Gravity},
Phys. Rev. D \textbf{112}, 054509 (2025).

\bibitem{LinearizedEFE2026}
J.~G.~Roehm,
\emph{Phase 6a Wave 1: Formal linearized Einstein equations and
emergent Newton constant from the ADW substrate} (in preparation).

\bibitem{Gilkey1995}
P.~B.~Gilkey,
\emph{Invariance Theory, the Heat Equation, and the Atiyah--Singer
Index Theorem}, 2nd ed.\ (CRC Press, 1995).

\bibitem{Vassilevich2003}
D.~V.~Vassilevich,
\emph{Heat kernel expansion: user's manual},
Phys. Rep. \textbf{388}, 279 (2003);
arXiv:hep-th/0306138.

\bibitem{ChristensenDuff1979}
S.~M.~Christensen and M.~J.~Duff,
\emph{New gravitational index theorems and supertheorems},
Nucl. Phys. B \textbf{154}, 301 (1979).

\end{thebibliography}

\end{document}
